\chapter{Weil sheaves and Weil complexes}\label{chap:weil-sheaves}

This chapter sets up the objects on which the Frobenius acts in the \(\ell\)-adic
proof of the Weil conjectures: the geometric fiber of a scheme over a finite
field, the Frobenius automorphism acting on it, and the categories of sheaves and
complexes equipped with a Frobenius structure isomorphism.

\section{The geometric fiber and its Galois Frobenius}

\begin{definition}[Geometric fiber]\label{def:weil-geometric-fiber}
\lean{AlgebraicGeometry.specAlgebraicClosure,
AlgebraicGeometry.Scheme.Hom.geometricFiber,
AlgebraicGeometry.Scheme.Hom.geometricFiberProj,
AlgebraicGeometry.Scheme.Hom.geometricFiberToAlgebraicClosure}
\leanok
Let \(k\) be a field with algebraic closure \(\bar k\) and let \(f\colon X_0\to
\operatorname{Spec}k\) be a scheme over \(k\).  The \emph{geometric fiber} of
\(f\) is the pullback
\[
 X:=X_0\times_{\operatorname{Spec}k}\operatorname{Spec}\bar k,
\]
with its projection \(p_X\colon X\to X_0\) to the model over \(k\) and its
structure morphism \(a_X\colon X\to\operatorname{Spec}\bar k\) over the algebraic
closure.
\end{definition}

\begin{definition}[Geometric Galois Frobenius]\label{def:weil-galois-frobenius}
\uses{def:weil-geometric-fiber}
\lean{FiniteField.geometricFrobenius,
AlgebraicGeometry.specGeometricFrobenius,
AlgebraicGeometry.Scheme.Hom.geometricGaloisFrobenius}
\leanok
Let \(k\) be a finite field.  The \emph{geometric Frobenius} of \(\bar k\) is the
inverse \(\varphi\) of the arithmetic Frobenius \(x\mapsto x^{\#k}\); it is a
\(k\)-algebra automorphism of \(\bar k\) and induces an automorphism
\(\sigma:=\operatorname{Spec}\varphi\) of \(\operatorname{Spec}\bar k\) over
\(\operatorname{Spec}k\).  The \emph{geometric Galois Frobenius} of the geometric
fiber of \(f\colon X_0\to\operatorname{Spec}k\) is its base change
\[
 g_X:=\mathbf 1_{X_0}\times\sigma\colon X\xrightarrow{\ \sim\ }X.
\]
It satisfies \(g_X\circ p_X=p_X\) and is semilinear over \(\bar k\), that is
\(a_X\circ g_X=\sigma\circ a_X\).
\end{definition}

\begin{remark}
The geometric Galois Frobenius is not the relative \(\#k\)-power Frobenius
\(\mathrm{Fr}_{X_0}\times\mathbf 1\) of \(X\) over \(\bar k\).  Since the
absolute Frobenius acts as the identity on the étale topos, the two induce
canonically isomorphic pullback endofunctors and hence equivalent categories of
Weil sheaves; the Galois version is used here because its Frobenius square is
Cartesian, which is what the compact-support base change needs.  The comparison
isomorphism, needed for the intersection-theoretic side of the Lefschetz trace
formula, is not yet formalized.
\end{remark}

\begin{lemma}[The Frobenius square is Cartesian]\label{lem:weil-frobenius-cartesian}
\uses{def:weil-galois-frobenius}
\lean{AlgebraicGeometry.Scheme.Hom.isPullback_geometricGaloisFrobenius}
\leanok
The square
\[
\begin{CD}
X @>{g_X}>> X\\
@V{a_X}VV @VV{a_X}V\\
\operatorname{Spec}\bar k @>{\sigma}>> \operatorname{Spec}\bar k
\end{CD}
\]
is Cartesian.
\end{lemma}
\begin{proof}
\leanok
The composite square obtained by pasting with the Cartesian square defining the
geometric fiber is Cartesian, because \(g_X\circ p_X=p_X\) and
\(\sigma\) is a morphism over \(\operatorname{Spec}k\); now cancel the right-hand
square.
\end{proof}

\section{Frobenius structures}

\begin{definition}[Algebra with invertible structure morphism]\label{def:weil-iso-algebra}
\lean{CategoryTheory.Endofunctor.Algebra.IsIsoStr,
CategoryTheory.Endofunctor.IsoAlgebra,
CategoryTheory.Endofunctor.IsoAlgebra.strIterate}
\leanok
Let \(F\colon\mathcal C\to\mathcal C\) be an endofunctor.  The category
\(\mathrm{IsoAlg}(F)\) is the full subcategory of the category of \(F\)-algebras
consisting of the algebras \((a,\Phi)\) whose structure morphism
\(\Phi\colon F(a)\to a\) is an isomorphism.  The \emph{iterated structure
isomorphism} is the composite
\(\Phi^{[n]}\colon F^n(a)\xrightarrow{\ \sim\ }a\), defined by
\(\Phi^{[0]}=\mathbf 1\) and \(\Phi^{[n+1]}=\Phi\circ F(\Phi^{[n]})\).
\end{definition}

\begin{definition}[Weil sheaf]\label{def:weil-sheaf}
\uses{def:weil-iso-algebra,def:weil-galois-frobenius}
\lean{AlgebraicGeometry.WeilSheaf}
\leanok
Let \(X\) be a scheme with an automorphism \(e\colon X\xrightarrow{\sim}X\) and
let \(A\) be a coefficient category.  A \emph{Weil sheaf} for \((X,e)\) is a
sheaf \(\mathcal F\) on the small étale site of \(X\) with values in \(A\)
together with an isomorphism \(\Phi\colon e^*\mathcal F\xrightarrow{\sim}\mathcal
F\); morphisms are morphisms of sheaves commuting with the structure
isomorphisms.  For \(X\) the geometric fiber of \(X_0\) over a finite field and
\(e=g_X\), this is the classical notion of a Weil sheaf \emph{on \(X_0\)}.
\end{definition}

\begin{definition}[Weil complex]\label{def:weil-complex-derived}
\uses{def:weil-iso-algebra,lem:weil-constructible-pullback}
\lean{AlgebraicGeometry.WeilComplex,
AlgebraicGeometry.WeilComplexPlus,
AlgebraicGeometry.WeilComplexBoundedConstructible}
\leanok
With \((X,e)\) as above, a \emph{Weil complex} for \((X,e)\) is an object \(K\)
of the derived category of étale sheaves on \(X\) together with an isomorphism
\(\Phi\colon e^*K\xrightarrow{\sim}K\), where \(e^*\) acts on the derived
category by termwise application, the pullback on sheaves being exact.
Restricting the pullback to the bounded below and to the bounded constructible
derived categories gives in the same way the categories of \emph{bounded below}
and of \emph{bounded constructible} Weil complexes, together with the comparison
functors
\(D^b_c\to D^{+}\to D\) between them.
\end{definition}

\begin{lemma}[The pullback of a constructible sheaf is constructible]
\label{lem:weil-constructible-sheaf-pullback}
\lean{AlgebraicGeometry.Scheme.isConstructibleSheaf_upperStar}
For a morphism of schemes \(g\colon Z\to Y\) and a constructible sheaf
\(\mathcal F\) on the small étale site of \(Y\), the pullback \(g^*\mathcal F\) is
constructible \cite[Tag 03SB]{stacks-project}.
\end{lemma}
\begin{proof}
Pull the trivializing stratification of \(Y\) back along \(g\): the base change
of an immersion is an immersion and the preimages of the strata partition \(Z\)
into locally closed subsets.  On each stratum the restriction of \(g^*\mathcal F\)
is the pullback of the restriction of \(\mathcal F\), and the pullback of a finite
locally constant sheaf is finite locally constant.  This proof is not yet
formalized.
\end{proof}

\begin{lemma}[The derived pullback preserves bounded constructible complexes]
\label{lem:weil-constructible-pullback}
\uses{lem:weil-constructible-sheaf-pullback}
\lean{AlgebraicGeometry.Etale.DerivedCategory.boundedConstructible_upperStar_mapDerivedCategory_obj,
AlgebraicGeometry.Etale.DerivedCategoryBoundedConstructible.upperStar}
\leanok
For a morphism of schemes \(g\colon Z\to Y\), the derived pullback \(g^*\) sends
\(D^b_c(Y,A)\) into \(D^b_c(Z,A)\), hence restricts to a functor
\(D^b_c(Y,A)\to D^b_c(Z,A)\).
\end{lemma}
\begin{proof}
\leanok
The pullback on sheaves is exact, so the derived pullback is \(t\)-exact and
preserves boundedness.  Its homology is the pullback of the homology, which is
constructible by \Cref{lem:weil-constructible-sheaf-pullback}.
\end{proof}

\begin{definition}[Homology of a Weil complex]\label{def:weil-homology}
\uses{def:weil-sheaf,def:weil-complex-derived}
\lean{AlgebraicGeometry.WeilComplex.homologyFunctor,
AlgebraicGeometry.WeilComplexBoundedConstructible.homologyFunctor}
\leanok
Since \(e^*\) is \(t\)-exact, the \(n\)-th homology sheaf of a Weil complex
carries the Weil sheaf structure
\(e^*H^n(K)\cong H^n(e^*K)\xrightarrow{H^n(\Phi)}H^n(K)\).  This defines functors
from Weil complexes, and from bounded constructible Weil complexes, to Weil
sheaves.
\end{definition}
